% =====================================================================
%  FK_BODY — Zusatzstudie: Die Funktionskritik (Diagnose und Programm)
% =====================================================================

\section{Anlass: Warum diese Studie zwischen Band VII und VIII steht}

Band VII endete mit einer Synthese, die zwei Lager gleichzeitig provozieren
muss. Der Manosphere nimmt sie die Schuldigen (\enquote{Reproduktion einer
Struktur ist keine Urheberschaft}), dem feministischen Diskurs mutet sie eine
Symmetrie zu (\enquote{geteilte Unfreiheit}, die aggregierte weibliche Wahl als
Stabilisatorin der Statusordnung). Bevor Band VIII diese Synthese in die
Zukunft verlängert --- Substitute, Staat, Re-Traditionalisierung ---, muss die
Reihe eine Rechenschaftsfrage beantworten: \emph{Wo steht sie eigentlich im
Verhältnis zur mächtigsten lebenden Theorietradition der Geschlechterfrage,
dem intersektionalen Feminismus?} Kollidieren die Befunde, oder sind die Ziele
am Ende deckungsgleich?

Diese Zusatzstudie leistet das in drei Teilen. \emph{Teil I} vermisst das
Verhältnis: die intersektionale Theorie in ihrer stärksten Form, dann die
ehrlichen Kollisionen, dann die ebenso ehrlichen Deckungen. \emph{Teil II}
prüft zwei schärfere Thesen: dass die funktionale Rolle des Mannes im
intersektionalen Rahmen systematisch vernachlässigt wird, und dass dieser
Rahmen strukturell \emph{pro Frau} argumentiert --- als Anwalt, nicht als
neutraler Gutachter. \emph{Teil III} zieht die Konsequenz und entwirft die
fehlende Disziplin: die \emph{Funktionskritik} --- ihren Gegenstand, ihre
Axiome, ihre Methode, ihre Politik und ihre stärksten Einwände gegen sich
selbst.

\begin{merksatz}
Die Leitfrage dieser Studie ist nicht, ob der Feminismus \enquote{recht hat}.
Sie lautet: Welche Arbeitsteilung zwischen Anwaltschaft und Strukturanalyse
braucht eine Geschlechterforschung, die niemanden unsichtbar macht --- und wer
leistet den Teil, der derzeit unbesetzt ist?
\end{merksatz}

% =====================================================================
\part*{Teil I --- Die Vermessung}
\addcontentsline{toc}{part}{Teil I --- Die Vermessung}

\section{Intersektionalität in ihrer stärksten Form}

Wer eine Theorie prüfen will, muss sie zuerst so darstellen, dass ihre
Vertreterinnen die Darstellung unterschreiben würden.

Der Begriff entsteht 1989 in einem \emph{rechtswissenschaftlichen} Kontext ---
und das ist für diese Reihe, die von der Rechtsstellung herkommt, kein Zufall,
sondern die erste Brücke. Kimberlé Crenshaw analysiert in
\textit{Demarginalizing the Intersection} den Fall \textit{DeGraffenreid v.\
General Motors}: Schwarze Frauen klagen gegen Diskriminierung und verlieren,
weil das Gericht die Kategorien einzeln prüft --- GM stellt Frauen ein (weiße,
im Büro) und Schwarze ein (Männer, in der Fertigung), also liege weder
Geschlechts- noch Rassendiskriminierung vor. \belegt\ Die schwarze Frau fällt
durch das Raster, weil das Raster aus Einzelkategorien besteht. Crenshaws
Pointe: Diskriminierung wirkt an \emph{Schnittstellen}, und wer nur entlang
einzelner Achsen prüft, macht genau die Menschen unsichtbar, die auf mehreren
Achsen zugleich stehen. In \textit{Mapping the Margins} (1991) wird daraus ein
allgemeines Programm; Patricia Hill Collins systematisiert es zur
\textit{matrix of domination} --- Herrschaftsverhältnisse als verschränkte,
nicht addierte Struktur; bell hooks liefert mit \enquote{imperialist white
supremacist capitalist patriarchy} die Formel für die Verschränkung von
Kapitalismus, Rassismus und Patriarchat; die deutsche Rezeption (Knapp;
Winker/Degele mit ihrer Mehrebenenanalyse) übersetzt das in ein
sozialwissenschaftliches Instrumentarium. \belegt

Vier Kernprinzipien tragen den Rahmen: \emph{Verschränkung} statt Addition
(Kategorien multiplizieren sich, sie stapeln sich nicht); \emph{Positionalität}
(Wissen ist standortgebunden, die marginalisierte Position sieht, was die
dominante übersieht --- Standpunkttheorie nach Harding/Hartsock);
\emph{Machtasymmetrie als Ausgangspunkt} (analysiert wird Herrschaft, nicht
Differenz); und \emph{Strukturebene vor Individualebene} (nicht der einzelne
Rassist ist der Gegenstand, sondern die Ordnung, die Rassismus normal macht).

Es lohnt, die Verwandtschaft sofort zu notieren: Die Funktionslogik-These
dieser Reihe ist ebenfalls eine Strukturtheorie, ebenfalls anti-essentialistisch
(die Funktion ist historisch variabel, nicht biologisch verhängt), und ihr
Gründungsbefund ist ebenfalls ein \emph{Rasterversagen}: Menschen fallen durch
Kategorien, die einzeln geprüft werden. Die Frage ist, ob die Verwandtschaft
die Kollisionen überlebt.

\section{Die Kollisionen}

\subsection{Asymmetrie: Axiom oder Variable?}

Die tiefste Kollision ist keine der Befunde, sondern der Architektur. Für den
intersektionalen Rahmen ist die Machtasymmetrie zugunsten des Mannes
\emph{Ausgangspunkt} der Analyse --- das Patriarchat ist die Voraussetzung, unter
der Verschränkungen untersucht werden. Für die Funktionskritik dieser Reihe
ist die Asymmetrie eine \emph{empirische Variable}: Sie kann je nach Feld,
Kohorte und Lebensphase in beide Richtungen zeigen --- die patriarchale
Dividende ist real (Pay Gap, Care-Verteilung, Gewalt gegen Frauen,
Repräsentanz), \emph{und} die Funktionslast ist real (Arbeitstod, Suizid,
Strafmaß, Wehrpflicht, Marktunsichtbarkeit), und welche Seite in einem
gegebenen Feld überwiegt, entscheidet die Messung, nicht das Axiom. \plausibel\
Band VII lieferte das Musterbeispiel: Im Alterssegment des Partnermarkts kippt
die Marktmacht demografisch zu den Männern --- die Asymmetrie folgt der
Knappheit, nicht dem Geschlecht. Ein Rahmen, der die Richtung der Asymmetrie
voraussetzt, kann solche Kippungen schwer verarbeiten.

\subsection{Das Agency-Problem}

Band VII sagt: Die aggregierte weibliche Partnerwahl stabilisiert die
Statusordnung. Der Satz ist sorgfältig gegen Schuldzuweisung abgesichert
(Bourdieu, Kandiyoti: Vollstreckung ist keine Urheberschaft) --- aber er
schreibt Frauen kausale Wirksamkeit zu. Im intersektionalen Diskurs steht
genau diese Zuschreibung unter Victim-Blaming-Verdacht: Wer die Beherrschten
als Mitreproduzenten beschreibt, so der Einwand, verschiebt Verantwortung von
der Struktur auf ihre Opfer. \umstritten\ Die Ironie: Kandiyotis
\textit{patriarchal bargain} und Connells \textit{emphasized femininity} ---
beides kanonische feministische Konzepte --- sagen exakt dasselbe. Die Kollision
liegt weniger in der Theorie als in der Diskursnorm: Was Kandiyoti über Frauen
im Nahen Osten schreiben durfte, gilt, auf westliche Partnerwahl angewandt,
schnell als Grenzverletzung.

\subsection{Methode: Verhaltensdaten gegen Standpunkt}

Band VII ruht auf quantitativer Verhaltensempirie --- Match-Raten,
Scheidungsregister, Einkommensknicks. Die intersektionale Forschung arbeitet
überwiegend qualitativ und standpunktepistemologisch: Die Erfahrung der
Marginalisierten hat evidenziellen Vorrang. Beide Methoden haben ihr Recht ---
aber sie kollidieren, wo Selbstauskunft und Verhalten auseinanderlaufen (die
stated/revealed-Lücke aus Band VII, §\,3.4). Und sie kollidieren im
Falsifikationsanspruch: Eine Theorie, die jeden Gegenbefund als Ausdruck der
Herrschaft reinterpretieren kann (mildere Strafen für Frauen als
\enquote{wohlwollender Sexismus}, also ebenfalls Patriarchat), immunisiert
sich gegen Widerlegung --- der Popper-Einwand, den auch interne Kritikerinnen
(McCall, Davis: \enquote{Intersectionality as buzzword}) kennen. \umstritten\
Fairerweise: Der Vorwurf trifft die Anwendungspraxis häufiger als die
Gründungstexte.

\subsection{Das Biologie-Tabu}

Band VII hält die Kontroverse Buss vs.\ Eagly/Wood offen --- evolutionäre
Disposition oder Rollenanpassung, die Daten entscheiden es nicht. Im
intersektionalen Rahmen gilt schon das Offenhalten vielfach als Biologismus.
Die Reihe kann diese Tür nicht schließen, ohne ihre Methode zu verraten:
\emph{Offenhalten ist keine Parteinahme.} \plausibel

\section{Die Deckungen}

\subsection{Dasselbe Rasterversagen}

Die überraschendste Deckung liegt im Gründungsbefund selbst. Crenshaws
schwarze Klägerin fällt durchs Raster, weil Kategorien einzeln geprüft werden.
Band VI/VII fanden dieselbe Figur am anderen Ende: Der arme, unsichtbare Mann
fällt durch das Raster von Klassenanalyse \emph{und} Geschlechteranalyse ---
für die eine ist er Mann (privilegiert), für die andere ist Geschlecht nicht
ihr Thema. Klasse\,$\times$\,Geschlecht macht ihn doppelt unsichtbar: die
DeGraffenreid-Figur, gespiegelt. \plausibel\ Es gibt Forschung, die genau hier
arbeitet --- Connells \textit{marginalized masculinities}, Kimmels
\textit{Angry White Men}, die Deaths-of-Despair-Literatur (Case/Deaton) ---,
und sie ist teils feministisch grundiert. Die Deckung ist real: Beide Rahmen
bekämpfen Kategorien, die Menschen an Schnittstellen unsichtbar machen.

\subsection{Gemeinsame Kronzeugen, gemeinsamer Gegner}

bell hooks steht in beiden Lagern: \textit{The Will to Change} beschreibt das
Patriarchat als System, das Männern emotionale Ganzheit gegen Macht abkauft ---
die \enquote{geteilte Unfreiheit} von Band VII, zwanzig Jahre früher und aus
der Mitte des Feminismus. Eva Illouz liefert die Brücke zur Marktanalyse: Ihre
Kritik der kommodifizierten Intimität (\textit{Warum Liebe endet}) und die
Zahlsklaven-These aus Band VII beschreiben dieselbe Bewirtschaftung des
Begehrens durch den Plattformkapitalismus. Und Nancy Frasers
\textit{universal caregiver model} formuliert das Entkopplungsziel bereits
symmetrisch: Eine gerechte Ordnung macht \emph{beide} Geschlechter zu
Sorgenden und Erwerbenden --- sie löst die Frau aus der Care-Funktion
\emph{und} den Mann aus der Ernährerfunktion. \belegt\ Der gemeinsame Gegner
beider Rahmen heißt Funktionsreduktion: Menschen auf ihre Funktion
festzulegen und ihren Wert daraus abzuleiten.

\begin{merksatz}
Ziel-Deckung: Was der Feminismus für die Frau fordert --- nicht auf die
Reproduktions- und Sorgefunktion reduziert zu werden ---, fordert die
Funktionskritik spiegelbildlich für den Mann: nicht auf die Ernährer-, Status-
und Opferfunktion reduziert zu werden. Es ist \emph{dieselbe} Forderung,
zweimal ausgesprochen.
\end{merksatz}

\begin{gegen}
Die Deckung darf nicht zur Vereinnahmung werden. Dass Fraser, hooks und
Connell die Bausteine liefern, heißt nicht, dass sie die Band-VII-Synthese
unterschreiben würden. Der Unterschied bleibt: Für den intersektionalen
Rahmen ist die Befreiung des Mannes ein \emph{Nebeneffekt} der Überwindung des
Patriarchats; für die Funktionskritik ist sie ein \emph{gleichrangiges
Analyseziel}. Nebeneffekt und Gleichrang sind verschiedene
Forschungsprogramme --- auch wenn sie am selben Horizont enden.
\end{gegen}

% =====================================================================
\part*{Teil II --- Die Leerstelle}
\addcontentsline{toc}{part}{Teil II --- Die Leerstelle}

\section{These 1: Die vernachlässigte Männerfunktion}

Die erste These lautet: Die funktionale Rolle des Mannes --- seine Bindung an
Ernährer-, Schutz- und Risikofunktionen samt deren Lasten --- wird im
intersektionalen Rahmen systematisch \emph{unterbeforscht und
unterinstitutionalisiert}. Sie lässt sich auf drei Ebenen prüfen.

\subsection{Die wissenschaftliche Ebene}

Die Männlichkeitsforschung existiert --- aber als Randbesatz. An den weit über
zweihundert Gender-Professuren des deutschsprachigen Raums ist sie nur
vereinzelt denominiert; es gibt keine eigenständige institutionalisierte
Männerforschung mit Lehrstühlen, Studiengängen oder Fachgesellschaft von
vergleichbarem Gewicht. \plausibel\ Ihre eigenen Vertreter beschreiben die
Marginalität seit Jahrzehnten --- von Hollsteins These der \enquote{vergessenen
Männer} bis zu Connells nüchterner Feststellung, dass die masculinity studies
im Gender-Feld ein Nischendasein führen. Der Befund ist kein Skandal, aber er
ist ein Befund: Die Kategorie Mann wird überwiegend als Privilegien- und
Täterposition prozessiert, selten als Funktions- und Verwundbarkeitsposition.

\subsection{Die institutionelle Ebene}

Härter ist die Infrastruktur-Bilanz. Den über 7\,000 Frauenhausplätzen in
Deutschland stehen bundesweit nur wenige Dutzend Männerschutzplätze in gut
einem Dutzend Einrichtungen gegenüber --- bei einem männlichen Anteil an den
Opfern häuslicher Gewalt, der amtlich um ein Fünftel bis Viertel liegt.
\plausibel\ Es gibt regelmäßige Frauengesundheitsberichte, aber keine
verstetigte Männergesundheitsberichterstattung --- bei rund fünf Jahren
Lebenserwartungslücke, einer Suizidrate von etwa drei männlichen Toten auf
eine weibliche und über 90\,\% Männeranteil an den tödlichen Arbeitsunfällen.
\belegt\ International: UN Women existiert; ein Pendant existiert nicht.
Gleichstellungspolitik ist operativ weitgehend Frauenförderpolitik ---
legitim begründbar, aber als Asymmetrie dokumentierbar: MINT-Programme für
Mädchen sind Standard, Care- und Grundschul-Programme für Jungen (Männeranteil
dort teils unter 15\,\%) die Ausnahme.

\subsection{Die diskursive Ebene}

Warum fällt die Leerstelle so wenig auf? Die Psychologie liefert Mechanismen.
Seager und Barry beschreiben als \emph{gamma bias} eine doppelte Verzerrung
der öffentlichen Geschlechterwahrnehmung: Männliches Leid wird minimiert,
männliche Täterschaft maximiert --- und spiegelbildlich weibliches Leid
vergrößert, weibliche Täterschaft verkleinert. \umstritten\ (Das Konzept ist
jung und nicht unwidersprochen, aber es bündelt robuste Einzelbefunde.)
Dazu passen die Moral-Typecasting-Forschung (Gray/Wegner: Wahrnehmung sortiert
in Täter- und Opferrollen, und Männer werden bevorzugt als Agenten, Frauen als
Patienten kodiert) und Rudman/Goodwins Befund eines starken automatischen
weiblichen Eigengruppen-Bias, dem kein männliches Gegenstück entspricht.
\belegt\ Die Sprachpolitik dokumentiert die Asymmetrie: \emph{Femizid} ist als
Kategorie etabliert --- angemessen für die spezifische Tötung von Frauen als
Frauen ---, während für die rund 80\,\% männlichen Mordopfer weltweit (UNODC)
keine geschlechtliche Kategorie existiert; der \emph{Gender Pay Gap} ist
Allgemeinwissen, der \emph{Gender Death Gap} kein Begriff. Und die vielleicht
härteste Einzelzahl: Sonja Starrs Analyse der US-Bundesstrafjustiz fand bei
gleicher Tat und Vorgeschichte im Schnitt \textbf{63\,\%} längere Strafen für
Männer --- eine Disparität, die jede vergleichbare rassische Disparität
übertrifft und diskursiv nahezu unsichtbar ist. \belegt

\section{These 2: Die Anwaltsarchitektur}

Die zweite These erklärt die erste: Die Leerstelle ist kein Versehen, sondern
folgt aus der Architektur. Feminismus ist --- dem Namen und der Geschichte nach
--- \emph{Anwaltschaft}: die organisierte Vertretung der Interessen von Frauen
gegen eine historisch dokumentierte Benachteiligung. Intersektionalität ist
das Analysewerkzeug dieser Anwaltschaft. Ein Anwalt prüft nicht neutral, wem
die Asymmetrie gerade dient; er vertritt seine Mandantin. Das ist keine
Korruption, es ist die Jobbeschreibung.

Daraus folgt die Prognose, die Teil II bestätigt hat: Befunde zulasten von
Frauen werden prozessiert und institutionalisiert; Befunde zulasten von
Männern werden entweder nicht erhoben, oder --- die elegantere Variante ---
als Patriarchatsfolge \emph{reinterpretiert}: Die mildere Bestrafung von
Frauen wird zum \enquote{wohlwollenden Sexismus}, der eigentlich Frauen
infantilisiert; der männliche Arbeitstod zur Folge männlicher
Risikosozialisation, also des Patriarchats. Die Reinterpretationen sind oft
sogar \emph{plausibel} --- das Patriarchat schädigt Männer wirklich, das ist
die These dieser Reihe seit Band VI. Aber als Diskursfigur hat
\enquote{patriarchy hurts men too} eine Doppelnatur: Sie kann Analyse sein
(hooks) oder Abwehrformel --- die Einsortierung männlichen Leids als
Unterpunkt, der keine eigene Forschung, kein eigenes Budget und keine eigene
Institution braucht, weil seine Lösung ja ohnehin mitgeliefert wird.
\plausibel\ Eine Erklärung, die jeden Befund absorbiert, hat aufgehört zu
erklären.

\begin{gegen}
Jetzt die Gegenbelege in ihrer stärksten Form --- sie sind stark. \emph{Erstens:}
Die Anwaltschaft ist historisch verdient. Frauen waren über die dokumentierte
Rechtsgeschichte dieser Reihe (Bände I--III) die rechtlich Geminderten;
nachholende Gerechtigkeit erzeugt notwendig eine temporäre Asymmetrie der
Aufmerksamkeit, und wer sie \enquote{Einseitigkeit} nennt, ohne die Geschichte
mitzunennen, verzerrt. \emph{Zweitens:} Männer dominieren die Parlamente,
Regierungen und Chefetagen seit jeher --- wenn Männerpolitik fehlt, haben
mehrheitlich Männer sie nicht gemacht. Die Leerstelle als feministisches
Versäumnis zu buchen, verwechselt Zuständigkeiten. \emph{Drittens:} Der Rahmen
ist korrekturfähig, und die Korrektur hat begonnen: Norwegens
\textit{Mannsutvalget} (NOU 2024:8) --- die erste staatliche Männerkommission,
feministisch mitgetragen --- hat männliche Problemlagen systematisch erfasst;
Richard Reeves' \textit{Of Boys and Men} (2022) wurde auch in feministischen
Medien ernsthaft rezipiert; Väterpolitik (Partnermonate) \emph{ist} real
existierende Männerpolitik. Wo korrigiert wird, war die Leerstelle Bug, nicht
Feature. \emph{Viertens:} Ein Teil der Kritik trifft nicht die Wissenschaft,
sondern ihre Vermischung mit Bewegungspolitik --- Gender Studies als Fach und
Feminismus als Bewegung sind nicht identisch, und die Studie sollte den
Unterschied nicht einebnen.
\end{gegen}

\section{Das Urteil: Feature mit Bug-Anteilen}

Die redlichste Zusammenfassung: Die Leerstelle ist \emph{strukturell} (sie
folgt aus der Anwaltsarchitektur --- insofern Feature), \emph{historisch
verstärkt} (nachholende Gerechtigkeit --- insofern kontingent) und
\emph{korrigierbar} (Mannsutvalget, Reeves --- insofern Bug). Daraus folgt die
entscheidende Weichenstellung für Teil III: Wenn die Anwaltschaft legitim ist
und bleiben soll, dann ist der Fehler nicht, dass der Feminismus seine
Mandantin vertritt. Der Fehler ist, dass die Gegenstelle unbesetzt ist ---
dass es keine etablierte Disziplin gibt, die die Funktionsbindung
\emph{beider} Geschlechter symmetrisch, verhaltensempirisch und ohne
Anwaltsauftrag untersucht. Man verklagt keinen Anwalt dafür, Anwalt zu sein.
Man wundert sich, dass das Gericht seit fünfzig Jahren ohne Gutachter tagt.

\begin{merksatz}
Nicht der Feminismus ist das Problem. Das Problem ist ein Diskurs, der
Anwaltschaft für Strukturanalyse hält --- und deshalb nicht bemerkt, dass die
Strukturanalyse fehlt.
\end{merksatz}

% =====================================================================
\part*{Teil III --- Der Entwurf: Die Funktionskritik}
\addcontentsline{toc}{part}{Teil III --- Der Entwurf: Die Funktionskritik}

\section{Drei Optionen --- und warum es die dritte sein muss}

Aus der Diagnose folgen drei mögliche Konsequenzen.

\emph{Option 1: Korrektur.} Der intersektionale Rahmen wird von innen
repariert --- mehr Männlichkeitsforschung, mehr Mannsutvalget, mehr Reeves.
Das geschieht bereits und ist wertvoll. Aber es lässt die Architektur
unberührt: Solange die Asymmetrie Axiom bleibt, bleibt männliche
Funktionslast ein Unterkapitel der Patriarchatskritik --- prozessiert nur,
soweit sie sich als dessen Nebenfolge erzählen lässt.

\emph{Option 2: Erweiterung.} Die \emph{Funktionsposition} wird als weitere
intersektionale Achse eingezogen: Geschlecht\,$\times$\,Klasse\,$\times$\,%
Funktionsbindung. Theoretisch elegant --- die unsichtbare männliche Mitte wäre
dann eine Schnittstellenposition wie jede andere. Aber die Erweiterung erbt
das Agency-Tabu und das Falsifikationsproblem des Wirtssystems; und sie
verlangt dem Rahmen ab, seine Ausgangsasymmetrie zur Variablen zu machen ---
also genau das aufzugeben, was ihn als Anwaltschaft funktionsfähig hält.

\emph{Option 3: Neugründung als Systemkritik.} Eine eigene Disziplin, deren
Gegenstand nicht Männer und nicht Frauen sind, sondern \emph{die
Funktionsordnung selbst}: das System, das Wert, Rechte und Begehrbarkeit an
geschlechtlich sortierte Funktionen koppelt. Sie konkurriert nicht mit der
Anwaltschaft --- sie besetzt die leere Gutachterbank. Das ist der Vorschlag
dieser Programmschrift.

\section{Die Disziplin: Gegenstand und Axiome}

\textbf{Name:} \emph{Funktionskritik} (voll: funktionskritische
Geschlechtertheorie). Der Name sagt Programm und Grenze zugleich: kritisiert
wird die Funktions\emph{ordnung}, nicht ein Geschlecht, nicht eine Bewegung.

\textbf{Gegenstand:} die Kopplung von Rechtsstellung, sozialem Wert und
partnerschaftlicher Begehrbarkeit an geschlechtlich zugewiesene Funktionen ---
historisch (Bände I--VI), gegenwärtig (Band VII) und prospektiv (Band VIII):
ihre Träger (Recht, Begehren, Markt, Staat), ihre Wanderungen und die
Bedingungen ihrer Auflösung.

\textbf{Vier Axiome} --- bewusst wenige, jedes falsifizierbar:

\begin{enumerate}
\item \textbf{Messbarkeit:} Funktionsbindung ist empirisch messbar (Prämien,
Lasten, Selektionsdaten) und historisch variabel. \emph{Falsifiziert, wenn}
sich Funktionszuweisungen als invariant gegenüber ökonomisch-institutionellem
Wandel erweisen.
\item \textbf{Asymmetrie als Variable:} Richtung und Größe der
Geschlechterasymmetrie sind je Feld zu messen, nicht vorauszusetzen.
\emph{Falsifiziert, wenn} die Messung feldübergreifend stets dieselbe Richtung
ergibt --- dann wäre das Axiom des Anwaltsrahmens empirisch bestätigt und die
Funktionskritik dessen Umweg.
\item \textbf{Agency ohne Schuld:} Akteure reproduzieren Strukturen durch
freie Akte, ohne deren Urheber zu sein; Kausalzuschreibung ist keine
Schuldzuschreibung. (Bourdieu/Kandiyoti als Fundament.)
\item \textbf{Symmetrie der Methode, nicht der Befunde:} Beide Geschlechter
werden mit identischen Instrumenten untersucht; die Befunde dürfen beliebig
asymmetrisch ausfallen. Symmetrie ist Eigenschaft des Verfahrens, nie ein
erwartetes Ergebnis.
\end{enumerate}

\textbf{Kernbegriffe:} \emph{Funktionsprämie} (Rechte/Status als Gegenleistung
für Funktion), \emph{Funktionslast} (Risiken/Pflichten derselben Kopplung),
\emph{Funktionsdividende} (verallgemeinerte patriarchale Dividende: der
Nettoertrag einer Funktionsposition, vorzeichenoffen), \emph{Statusselektion}
(Partnerwahl als Vollzug der Funktionsordnung, Band VII),
\emph{geteilte Unfreiheit} (die Systemlage ohne reine Gewinner) und das
\emph{Entkopplungsziel}: eine Ordnung, in der Wert und Funktion getrennt sind
--- Frasers universal caregiver, symmetrisch ergänzt um den universal
provider.

\section{Der Methodenkanon}

Die Methode der Reihe, als Disziplinregeln kodifiziert:

\begin{enumerate}
\item \textbf{Verhalten vor Selbstauskunft.} Wo geäußerte und gelebte
Präferenzen divergieren, entscheidet das Verhalten (Registerdaten,
Marktdaten); Selbstauskünfte werden als Signale analysiert, nicht als Messwerte
übernommen.
\item \textbf{Konfidenzpflicht.} Jede Aussage trägt ihre Stufe --- belegt,
plausibel, umstritten, spekulativ. Unmarkierte Gewissheit ist ein
Methodenfehler.
\item \textbf{Gegendarstellungspflicht.} Jede tragende These wird von ihrem
stärksten Einwand begleitet (Steelmanning); ein Text ohne Gegendarstellung
gilt als unfertig.
\item \textbf{Falsifikationsangebot.} Jede Kernthese benennt, welcher Befund
sie widerlegen würde. Was nichts ausschließt, behauptet nichts.
\item \textbf{Schuldrhetorik-Verbot.} Strukturbefunde werden nie als
Kollektivschuld formuliert --- weder gegen Frauen (Band VII) noch gegen Männer
noch gegen Bewegungen.
\item \textbf{Kooperationsgebot.} Die Funktionskritik ersetzt keine
Anwaltschaft und bekämpft keine; sie liefert die Strukturanalyse, auf die
Anwaltschaften beider Richtungen sich beziehen können. Arbeitsteilung, nicht
Konkurrenz.
\end{enumerate}

\section{Forschungsagenda: zehn Prüffragen}

\begin{enumerate}
\item Sinkt die weibliche Statuspräferenz messbar, wo Pay Gap und Care-Risiko
sinken (natürliche Experimente Skandinavien)? --- Der direkte Test von Axiom 1
gegen die Disposition-These.
\item Repliziert sich Starrs Strafzumessungs-Gap in deutschen Registerdaten,
und wie verändert er sich im Zeitverlauf?
\item Lässt sich der gamma bias experimentell in deutschen Stichproben
zeigen --- und ist er durch Information korrigierbar?
\item Wie hoch sind die realen Karrieren- und Partnerwertkosten männlicher
Elternzeit über zwei Monate (Längsschnitt)?
\item Kippt die Partnermarkt-Asymmetrie im Alterssegment wirklich (Axiom 2 am
Band-VII-Befund), und wie verhält sich dann das Commitment-Verhalten?
\item Welche Wirkungen zeigt das norwegische Mannsutvalget nach fünf Jahren ---
und was davon ist auf Deutschland übertragbar?
\item Begriffsfeldanalyse: Wie strukturiert die Asymmetrie Femizid/kein
Pendant, Pay Gap/kein Death Gap die Medienberichterstattung messbar?
\item Ist die KI-Companion-Nutzung funktionsflüchtig (Ausweichen vor
Statusselektion) oder funktionsäquivalent (Substitution der Funktion) ---
Nutzerdaten nach Marktposition? (Vorgriff auf Band VIII.)
\item Wirkt die neue Wehrerfassung 2026ff.\ als Funktionsrückkehr --- messbar
in Einstellungen junger Männer zu Staat und Geschlechterordnung? (Band VIII.)
\item Entkopplungs-Indikatorik: Welche messbaren Größen zeigen an, dass Wert
und Funktion sich trennen (z.\,B. Trennungsrisiko bei Ernährerwechsel,
Partnerwert von Care-Männern) --- der Fortschrittsmaßstab der Disziplin.
\end{enumerate}

\section{Politische Übersetzung}

Die Funktionskritik ist Analyse, keine Partei --- aber sie hat einen
Prüfmaßstab für Politik: \emph{Entkoppelt die Maßnahme Wert von Funktion, für
beide Geschlechter?} Daraus folgen symmetrische Prüfsteine: Frasers universal
caregiver ernst nehmen heißt Partnermonate ausbauen \emph{und} die
Care-Berufe für Männer öffnen (Jungen-Programme dort, wo der Männeranteil
unter 15\,\% liegt); Schutzinfrastruktur nach Bedarf statt nach Kategorie
(Frauenhausplätze ausbauen \emph{und} Männerschutzplätze aus der
Alibigröße holen); verstetigte Männergesundheitsberichterstattung neben der
weiblichen; und der Testfall des Jahrzehnts: die Wehrpflichtfrage --- eine
Ordnung, die die letzte exklusive Geschlechtspflicht behält, hat die
Entkopplung nicht vollzogen, egal wie sie sich nennt. \plausibel\ Norwegens
Mannsutvalget zeigt das Institutionenmodell: keine Gegenbewegung, sondern eine
staatliche Kommission, feministisch mitgetragen, symmetrisch beauftragt.

\section{Selbstkritik: die fünf stärksten Einwände}

\begin{gegen}
\textbf{1 · Verfrühte Symmetrie.} Solange die patriarchale Dividende real ist,
legitimiert methodische Symmetrie faktische Asymmetrie --- der
All-Lives-Matter-Einwand. \emph{Antwort:} Axiom 4 erwartet asymmetrische
Befunde ausdrücklich; die Dividende ist Bestandteil der Bilanz, nicht ihr
Gegner. Aber der Einwand bleibt als Wachsamkeitsregel gültig: Symmetrie der
Methode darf nie als Behauptung gleicher Lasten verkauft werden.

\textbf{2 · Vereinnahmung von rechts.} Jede Männer-Leerstellen-Diagnose wird
von Maskulisten als Munition gelesen werden. \emph{Antwort:} Die
Schutzmechanismen sind eingebaut --- Schuldrhetorik-Verbot, Dividenden-Anerkennung,
Kooperationsgebot --- und die Erfahrung von Band VII zeigt: Die beste
Verteidigung gegen Vereinnahmung ist die Gegendarstellungspflicht im Text
selbst. Wer die Reihe zitiert, zitiert die Einwände mit.

\textbf{3 · Die Funktionalismus-Falle.} Wer alles als Funktion erklärt, lässt
die Ordnung alternativlos erscheinen --- Funktionalismus war schon einmal die
Ideologie des Status quo. \emph{Antwort:} Axiom 1 (historische Variabilität)
und das Entkopplungsziel sind die Gegengifte; die Disziplin misst
Funktionsbindung, um sie auflösbar zu machen, nicht um sie zu rechtfertigen.
Der naturalistische Fehlschluss --- aus \emph{ist} folgt \emph{soll} --- wird
in Band VIII eigens behandelt.

\textbf{4 · Empirismus-Blindheit.} Verhaltensdaten sehen keine Erfahrung; wer
nur Register liest, verliert, was Standpunkttheorie sichtbar macht.
\emph{Antwort:} berechtigt --- deshalb die dramatisierten Szenen der Reihe als
methodisches Gegengewicht, und deshalb Kooperations- statt
Ersetzungsanspruch gegenüber qualitativer Forschung.

\textbf{5 · Akademische Hybris.} Braucht es eine \enquote{Disziplin}, oder
reicht Reeves' Pragmatik --- Probleme benennen, Institutionen bauen, ohne
Theoriegebäude? \emph{Antwort:} Vielleicht reicht sie. Der Entwurf ist ein
Angebot, kein Anspruch; wenn seine Prüffragen von bestehenden Fächern
abgearbeitet werden, hat er seinen Zweck erfüllt, ohne je gegründet worden zu
sein.
\end{gegen}

\section{Schluss: Verhältnisbestimmung und Übergabe}

Die Antwort auf die Ausgangsfrage lautet also: Die Band-VII-Synthese
\emph{kollidiert} mit dem intersektionalen Feminismus in der Architektur
(Asymmetrie-Axiom, Agency-Tabu, Standpunkt-Methode) --- und ist mit ihm
\emph{zieldeckungsgleich} im Kern (Anti-Essentialismus, Kampf gegen
Rasterversagen, Entkopplung von Wert und Funktion). Die Konsequenz ist weder
Unterwerfung noch Gegnerschaft, sondern Arbeitsteilung: Der Feminismus bleibt
die legitime Anwaltschaft der historisch Geminderten; die Funktionskritik
besetzt die leere Gutachterbank --- als Systemkritik, deren Gegenstand die
Funktionsordnung selbst ist und deren Erfolgsmaß die Entkopplung für
\emph{alle} ist.

\begin{merksatz}
Die Funktionskritik ist kein Anti-Feminismus. Sie ist das zweite Auge einer
Geschlechterforschung, die bisher mit einem sah --- notwendig geworden nicht,
weil das erste Auge schlecht sähe, sondern weil Tiefenschärfe zwei braucht.
\end{merksatz}

Damit ist die Reihe für Band VIII gerüstet: Wenn die Funktionsordnung ihre
Träger wechselt --- vom Recht über das Begehren zu Substituten, Staat und
re-inszenierter Tradition ---, dann hat sie jetzt die Disziplin, die diese
Wanderung messen kann.

\section*{Quellen und Konfidenz}
\addcontentsline{toc}{section}{Quellen und Konfidenz}
\small
Tragende Belege dieser Zusatzstudie; Konfidenz jeweils markiert. Vor
Drucklegung sind die institutionellen Zahlen (Schutzplätze, Professuren) gegen
die aktuellen Primärquellen zu verifizieren.

\begin{itemize}
\item \textbf{Crenshaw, K. (1989/1991):} \textit{Demarginalizing the
Intersection} / \textit{Mapping the Margins}; DeGraffenreid v.\ GM. \belegt.
\item \textbf{Collins, P.\,H.:} matrix of domination; \textbf{hooks, b.:}
\textit{The Will to Change} (2004); \textbf{Winker/Degele:}
Mehrebenenanalyse. \belegt\ (Theorie).
\item \textbf{Kandiyoti, D. (1988); Bourdieu, P. (1998); Connell, R.:}
patriarchal bargain, symbolische Gewalt, hegemoniale/marginalisierte
Männlichkeit, patriarchale Dividende. \belegt\ (Theorie).
\item \textbf{Fraser, N.:} universal caregiver model; \textbf{Illouz, E.:}
Kommodifizierung der Intimität. \belegt\ (Theorie).
\item \textbf{Starr, S. (2012/2015):} Gender Disparities in Federal Criminal
Cases --- \textasciitilde 63\,\% längere Strafen für Männer. \belegt.
\item \textbf{UNODC Global Study on Homicide:} \textasciitilde 80\,\% der
Mordopfer männlich. \belegt.
\item \textbf{Destatis/DGUV:} Suizidverhältnis \textasciitilde 3:1;
Arbeitsunfalltote $>$90\,\% männlich; Lebenserwartungslücke
\textasciitilde 5 Jahre. \belegt.
\item \textbf{Schutzinfrastruktur:} $>$7\,000 Frauenhausplätze vs.\ wenige
Dutzend Männerschutzplätze. \plausibel\ (Größenordnung; Zählungen variieren).
\item \textbf{Seager, M. \& Barry, J.:} gamma bias, Palgrave Handbook of Male
Psychology. \umstritten\ (junges Konzept).
\item \textbf{Gray/Wegner:} moral typecasting; \textbf{Rudman/Goodwin:}
weiblicher In-group-Bias. \belegt.
\item \textbf{Benatar, D. (2012):} \textit{The Second Sexism} --- samt der
Kritik daran. \umstritten.
\item \textbf{NOU 2024:8 (Mannsutvalget):} Norwegens Männerkommission.
\belegt.
\item \textbf{Reeves, R. (2022):} \textit{Of Boys and Men}; AIBM. \belegt.
\item \textbf{McCall, L.; Davis, K.:} interne Intersektionalitätskritik
(Komplexität, buzzword). \belegt\ (Debattenlage).
\end{itemize}

\vspace{1em}
{\color{obsmuted}\rule{\textwidth}{0.4pt}}\\[0.4em]
{\ttfamily\scriptsize\color{obsmuted}
OBSERVATORY · Zusatzstudie zu Band VII \enquote{Die Funktionskritik} ·
Diagnose und Programm · Leitformel: Rechtsstellung folgt Funktionsstellung ·
Diese Programmschrift ist ein Angebot zur Arbeitsteilung, kein Angriff auf
eine Bewegung; jede ihrer Thesen trägt Konfidenzmarker und Gegendarstellung.}
